% ============================================
% GIFT v3.5 — Shared Bibliography
% ============================================
% Usage: % ============================================
% GIFT v3.5 — Shared Bibliography
% ============================================
% Usage: % ============================================
% GIFT v3.5 — Shared Bibliography
% ============================================
% Usage: % ============================================
% GIFT v3.5 — Shared Bibliography
% ============================================
% Usage: \input{gift_3.5_bib} where the references section should appear.
% Sequential numbering [1]..[50]. Round-1 + round-2 IA-review updates applied
% (2026-07-07): [13] Kovalev 2005, [26] T2K/NOvA Nature 646,
% [27] NuFit 6.0 + 6.1, [45] Garbagnati-Sarti 2007, [46] LEPSUSYWG,
% [47] Fermilab DUNE 2031, [48] Chen-Hong (re-added, cited in S1), [49] Planck PR4, [50] CODATA.

\begin{thebibliography}{99}

\bibitem{1} Particle Data Group, ``Review of Particle Physics,'' \textit{Phys.\ Rev.\ D} \textbf{110}, 030001 (2024).

\bibitem{2} S.~Weinberg, ``Implications of dynamical symmetry breaking,'' \textit{Phys.\ Rev.\ D} \textbf{13}, 974 (1976).

\bibitem{3} Planck Collaboration, ``Planck 2018 results.\ VI.\ Cosmological parameters,'' \textit{Astron.\ Astrophys.} \textbf{641}, A6 (2020).

\bibitem{4} A.~G.~Riess et al., ``A comprehensive measurement of the local value of the Hubble constant,'' \textit{ApJL} \textbf{934}, L7 (2022).

\bibitem{5} C.~D.~Froggatt, H.~B.~Nielsen, ``Hierarchy of quark masses, Cabibbo angles and CP violation,'' \textit{Nucl.\ Phys.\ B} \textbf{147}, 277 (1979).

\bibitem{6} Y.~Koide, ``Fermion-boson two-body model of quarks and leptons and Cabibbo mixing,'' \textit{Lett.\ Nuovo Cim.} \textbf{34}, 201 (1982).

\bibitem{7} C.~Furey, \textit{Standard Model Physics from an Algebra?}, PhD thesis, University of Waterloo (2015).

\bibitem{8} N.~Furey, M.~J.~Hughes, ``One generation of standard model Weyl representations as a single copy of $\mathbb{R} \otimes \mathbb{C} \otimes \mathbb{H} \otimes \mathbb{O}$,'' \textit{Phys.\ Lett.\ B} \textbf{831}, 137186 (2022).

\bibitem{9} R.~Wilson, ``$E_8$ and Standard Model plus gravity,'' arXiv:2404.18938 [hep-th] (2024).

\bibitem{10} T.~P.~Singh et al., ``An $E_8 \otimes E_8$ unification of the Standard Model with pre-gravitation,'' arXiv:2206.06911v3 (2024).

\bibitem{11} B.~S.~Acharya, S.~Gukov, ``M-theory and singularities of exceptional holonomy manifolds,'' \textit{Phys.\ Rep.} \textbf{392}, 121 (2004).

\bibitem{12} L.~Foscolo et al., ``Complete non-compact $G_2$-manifolds from asymptotically conical Calabi--Yau 3-folds,'' \textit{Duke Math.\ J.} \textbf{170}, 3 (2021).

\bibitem{13} A.~Kovalev, ``Coassociative K3 fibrations of compact $G_2$-manifolds,'' arXiv:math/0511150 (2005).

\bibitem{14} M.~Haskins, J.~Nordstr\"om, ``Extra-twisted connected sum $G_2$-manifolds,'' arXiv:1809.09083 (2022); \textit{Ann.\ Glob.\ Anal.\ Geom.} \textbf{64}, art.\ 2 (2023).

\bibitem{15} G.~Bera, ``Associative submanifolds in twisted connected sum $G_2$-manifolds,'' arXiv:2209.00156 (2022).

\bibitem{16} D.~Crowley, S.~Goette, J.~Nordstr\"om, ``An analytic invariant of $G_2$ manifolds,'' \textit{Invent.\ Math.} \textbf{239}(3), 865--907 (2025).

\bibitem{17} T.~Dray, C.~A.~Manogue, \textit{The Geometry of the Octonions}, World Scientific (2015).

\bibitem{18} J.~F.~Adams, \textit{Lectures on Exceptional Lie Groups}, University of Chicago Press (1996).

\bibitem{19} D.~J.~Gross et al., ``Heterotic string theory,'' \textit{Nucl.\ Phys.\ B} \textbf{256}, 253 (1985).

\bibitem{20} D.~D.~Joyce, \textit{Compact Manifolds with Special Holonomy}, Oxford University Press (2000).

\bibitem{21} A.~Kovalev, ``Twisted connected sums and special Riemannian holonomy,'' \textit{J.\ Reine Angew.\ Math.} \textbf{565}, 125 (2003).

\bibitem{22} A.~Corti et al., ``$G_2$-manifolds and associative submanifolds via semi-Fano 3-folds,'' \textit{Duke Math.\ J.} \textbf{164}, 1971 (2015).

\bibitem{23} R.~Harvey, H.~B.~Lawson, ``Calibrated geometries,'' \textit{Acta Math.} \textbf{148}, 47 (1982).

\bibitem{24} B.~S.~Acharya, ``M-theory, $G_2$-manifolds and four-dimensional physics,'' \textit{Class.\ Quant.\ Grav.} \textbf{19}, 5619 (2002).

\bibitem{25} B.~S.~Acharya, E.~Witten, ``Chiral fermions from manifolds of $G_2$ holonomy,'' arXiv:hep-th/0109152 (2001).

\bibitem{26} T2K and NOvA Collaborations, ``Joint analysis of long-baseline neutrino oscillations,'' \textit{Nature} \textbf{646}, 818--824 (2025), \href{https://doi.org/10.1038/s41586-025-09599-3}{doi:10.1038/s41586-025-09599-3}.

\bibitem{27} I.~Esteban et al.\ (NuFIT), ``NuFit-6.0: three-flavour global analysis of neutrino oscillation experiments,'' \textit{JHEP} \textbf{12} (2024) 216, arXiv:2410.05380; NuFIT 6.1 web release, \url{https://www.nu-fit.org} (2025).

\bibitem{28} L.~de Moura, S.~Ullrich, ``The Lean 4 theorem prover and programming language,'' in \textit{CADE 28}, p.\ 625 (2021).

\bibitem{29} mathlib Community, \textit{mathlib4}, \url{https://github.com/leanprover-community/mathlib4}.

\bibitem{30} DUNE Collaboration, ``Long-baseline neutrino facility (LBNF) and DUNE conceptual design report,'' FERMILAB-TM-2696 (2020).

\bibitem{31} DUNE Collaboration, ``Deep Underground Neutrino Experiment (DUNE) far detector technical design report,'' arXiv:2103.04797 (2021).

\bibitem{32} E.~Witten, ``Strong coupling expansion of Calabi-Yau compactification,'' \textit{Nucl.\ Phys.\ B} \textbf{471}, 135 (1996).

\bibitem{33} B.~S.~Acharya et al., ``M-theory solution to the hierarchy problem,'' \textit{Phys.\ Rev.\ D} \textbf{76}, 126010 (2007).

\bibitem{34} M.~Atiyah, E.~Witten, ``M-theory dynamics on a manifold of $G_2$-holonomy,'' \textit{Adv.\ Theor.\ Math.\ Phys.} \textbf{6}, 1 (2002).

\bibitem{35} G.~Kane, \textit{String Theory and the Real World} (2017).

\bibitem{36} J.~Distler, S.~Garibaldi, ``There is no `Theory of Everything' inside $E_8$,'' \textit{Commun.\ Math.\ Phys.} \textbf{298}, 419 (2010).

\bibitem{37} J.~C.~Baez, ``Octonions and the Standard Model,'' \url{https://math.ucr.edu/home/baez/standard/} (2020--2025).

\bibitem{38} Springer Nature, ``Artificial intelligence (AI) policy,'' \url{https://www.springernature.com/gp/policies} (2024).

\bibitem{39} J.~A.~Wheeler, ``Information, physics, quantum: the search for links,'' in \textit{Complexity, Entropy, and the Physics of Information}, W.~H.~Zurek (ed.), Addison-Wesley (1990), pp.~3--28.

\bibitem{40} J.~Worrall, ``Structural realism: the best of both worlds?,'' \textit{Dialectica} \textbf{43}, 99--124 (1989).

\bibitem{41} J.~Ladyman, D.~Ross, \textit{Every Thing Must Go: Metaphysics Naturalized}, Oxford University Press (2007).

\bibitem{42} R.~Pin\v{c}\'ak, A.~Pigazzini, M.~Pudl\'ak, E.~Barto\v{s}, ``Geometric origin of a stable black hole remnant from torsion in $G_2$-manifold geometry,'' \textit{Gen.\ Rel.\ Grav.} \textbf{58}, 29 (2026), \href{https://doi.org/10.1007/s10714-026-03528-z}{doi:10.1007/s10714-026-03528-z}.

\bibitem{43} D.~Joyce, S.~Karigiannis, ``A new construction of compact torsion-free $G_2$-manifolds by gluing families of Eguchi--Hanson spaces,'' \textit{J.\ Differential Geom.} (2021), arXiv:1707.09325 (2017).

\bibitem{44} S.~Mukai, ``Finite groups of automorphisms of K3 surfaces and the Mathieu group,'' \textit{Invent.\ Math.} \textbf{94}, 183--221 (1988).

\bibitem{45} A.~Garbagnati, A.~Sarti, ``Symplectic automorphisms of prime order on K3 surfaces,'' \textit{J.\ Algebra} \textbf{318}, 323--350 (2007), arXiv:math/0603742.

\bibitem{46} LEP SUSY Working Group (ALEPH, DELPHI, L3, OPAL), ``Combined LEP chargino results, up to 208 GeV for large $m_0$,'' LEPSUSYWG/01-03.1 (2001), \url{https://lepsusy.web.cern.ch/lepsusy/www/inos_moriond01/charginos_pub.html}.

\bibitem{47} Fermilab, ``DUNE first-beam target of 2031: LBNF/DUNE-US milestone communication'' (2026-05-07), \url{https://news.fnal.gov/}.

\bibitem{48} J.~Chen, H.~Hong, ``Intermediate curvature and splitting theorem,'' arXiv:2604.26529 (2026).

\bibitem{49} M.~Tristram et al., ``Cosmological parameters derived from the final Planck data release (PR4),'' \textit{Astron.\ Astrophys.} \textbf{682}, A37 (2024).

\bibitem{50} E.~Tiesinga, P.~J.~Mohr, D.~B.~Newell, B.~N.~Taylor, ``CODATA recommended values of the fundamental physical constants: 2022,'' \textit{Rev.\ Mod.\ Phys.} \textbf{97}, 025002 (2025), \url{https://physics.nist.gov/cuu/Constants}.

\bibitem{51} V.~V.~Nikulin, ``Integer symmetric bilinear forms and some of their geometric applications,'' \textit{Izv.\ Akad.\ Nauk SSSR Ser.\ Mat.} \textbf{43} (1979); English transl.\ \textit{Math.\ USSR Izv.} \textbf{14} (1980).

\end{thebibliography}
 where the references section should appear.
% Sequential numbering [1]..[50]. Round-1 + round-2 IA-review updates applied
% (2026-07-07): [13] Kovalev 2005, [26] T2K/NOvA Nature 646,
% [27] NuFit 6.0 + 6.1, [45] Garbagnati-Sarti 2007, [46] LEPSUSYWG,
% [47] Fermilab DUNE 2031, [48] Chen-Hong (re-added, cited in S1), [49] Planck PR4, [50] CODATA.

\begin{thebibliography}{99}

\bibitem{1} Particle Data Group, ``Review of Particle Physics,'' \textit{Phys.\ Rev.\ D} \textbf{110}, 030001 (2024).

\bibitem{2} S.~Weinberg, ``Implications of dynamical symmetry breaking,'' \textit{Phys.\ Rev.\ D} \textbf{13}, 974 (1976).

\bibitem{3} Planck Collaboration, ``Planck 2018 results.\ VI.\ Cosmological parameters,'' \textit{Astron.\ Astrophys.} \textbf{641}, A6 (2020).

\bibitem{4} A.~G.~Riess et al., ``A comprehensive measurement of the local value of the Hubble constant,'' \textit{ApJL} \textbf{934}, L7 (2022).

\bibitem{5} C.~D.~Froggatt, H.~B.~Nielsen, ``Hierarchy of quark masses, Cabibbo angles and CP violation,'' \textit{Nucl.\ Phys.\ B} \textbf{147}, 277 (1979).

\bibitem{6} Y.~Koide, ``Fermion-boson two-body model of quarks and leptons and Cabibbo mixing,'' \textit{Lett.\ Nuovo Cim.} \textbf{34}, 201 (1982).

\bibitem{7} C.~Furey, \textit{Standard Model Physics from an Algebra?}, PhD thesis, University of Waterloo (2015).

\bibitem{8} N.~Furey, M.~J.~Hughes, ``One generation of standard model Weyl representations as a single copy of $\mathbb{R} \otimes \mathbb{C} \otimes \mathbb{H} \otimes \mathbb{O}$,'' \textit{Phys.\ Lett.\ B} \textbf{831}, 137186 (2022).

\bibitem{9} R.~Wilson, ``$E_8$ and Standard Model plus gravity,'' arXiv:2404.18938 [hep-th] (2024).

\bibitem{10} T.~P.~Singh et al., ``An $E_8 \otimes E_8$ unification of the Standard Model with pre-gravitation,'' arXiv:2206.06911v3 (2024).

\bibitem{11} B.~S.~Acharya, S.~Gukov, ``M-theory and singularities of exceptional holonomy manifolds,'' \textit{Phys.\ Rep.} \textbf{392}, 121 (2004).

\bibitem{12} L.~Foscolo et al., ``Complete non-compact $G_2$-manifolds from asymptotically conical Calabi--Yau 3-folds,'' \textit{Duke Math.\ J.} \textbf{170}, 3 (2021).

\bibitem{13} A.~Kovalev, ``Coassociative K3 fibrations of compact $G_2$-manifolds,'' arXiv:math/0511150 (2005).

\bibitem{14} M.~Haskins, J.~Nordstr\"om, ``Extra-twisted connected sum $G_2$-manifolds,'' arXiv:1809.09083 (2022); \textit{Ann.\ Glob.\ Anal.\ Geom.} \textbf{64}, art.\ 2 (2023).

\bibitem{15} G.~Bera, ``Associative submanifolds in twisted connected sum $G_2$-manifolds,'' arXiv:2209.00156 (2022).

\bibitem{16} D.~Crowley, S.~Goette, J.~Nordstr\"om, ``An analytic invariant of $G_2$ manifolds,'' \textit{Invent.\ Math.} \textbf{239}(3), 865--907 (2025).

\bibitem{17} T.~Dray, C.~A.~Manogue, \textit{The Geometry of the Octonions}, World Scientific (2015).

\bibitem{18} J.~F.~Adams, \textit{Lectures on Exceptional Lie Groups}, University of Chicago Press (1996).

\bibitem{19} D.~J.~Gross et al., ``Heterotic string theory,'' \textit{Nucl.\ Phys.\ B} \textbf{256}, 253 (1985).

\bibitem{20} D.~D.~Joyce, \textit{Compact Manifolds with Special Holonomy}, Oxford University Press (2000).

\bibitem{21} A.~Kovalev, ``Twisted connected sums and special Riemannian holonomy,'' \textit{J.\ Reine Angew.\ Math.} \textbf{565}, 125 (2003).

\bibitem{22} A.~Corti et al., ``$G_2$-manifolds and associative submanifolds via semi-Fano 3-folds,'' \textit{Duke Math.\ J.} \textbf{164}, 1971 (2015).

\bibitem{23} R.~Harvey, H.~B.~Lawson, ``Calibrated geometries,'' \textit{Acta Math.} \textbf{148}, 47 (1982).

\bibitem{24} B.~S.~Acharya, ``M-theory, $G_2$-manifolds and four-dimensional physics,'' \textit{Class.\ Quant.\ Grav.} \textbf{19}, 5619 (2002).

\bibitem{25} B.~S.~Acharya, E.~Witten, ``Chiral fermions from manifolds of $G_2$ holonomy,'' arXiv:hep-th/0109152 (2001).

\bibitem{26} T2K and NOvA Collaborations, ``Joint analysis of long-baseline neutrino oscillations,'' \textit{Nature} \textbf{646}, 818--824 (2025), \href{https://doi.org/10.1038/s41586-025-09599-3}{doi:10.1038/s41586-025-09599-3}.

\bibitem{27} I.~Esteban et al.\ (NuFIT), ``NuFit-6.0: three-flavour global analysis of neutrino oscillation experiments,'' \textit{JHEP} \textbf{12} (2024) 216, arXiv:2410.05380; NuFIT 6.1 web release, \url{https://www.nu-fit.org} (2025).

\bibitem{28} L.~de Moura, S.~Ullrich, ``The Lean 4 theorem prover and programming language,'' in \textit{CADE 28}, p.\ 625 (2021).

\bibitem{29} mathlib Community, \textit{mathlib4}, \url{https://github.com/leanprover-community/mathlib4}.

\bibitem{30} DUNE Collaboration, ``Long-baseline neutrino facility (LBNF) and DUNE conceptual design report,'' FERMILAB-TM-2696 (2020).

\bibitem{31} DUNE Collaboration, ``Deep Underground Neutrino Experiment (DUNE) far detector technical design report,'' arXiv:2103.04797 (2021).

\bibitem{32} E.~Witten, ``Strong coupling expansion of Calabi-Yau compactification,'' \textit{Nucl.\ Phys.\ B} \textbf{471}, 135 (1996).

\bibitem{33} B.~S.~Acharya et al., ``M-theory solution to the hierarchy problem,'' \textit{Phys.\ Rev.\ D} \textbf{76}, 126010 (2007).

\bibitem{34} M.~Atiyah, E.~Witten, ``M-theory dynamics on a manifold of $G_2$-holonomy,'' \textit{Adv.\ Theor.\ Math.\ Phys.} \textbf{6}, 1 (2002).

\bibitem{35} G.~Kane, \textit{String Theory and the Real World} (2017).

\bibitem{36} J.~Distler, S.~Garibaldi, ``There is no `Theory of Everything' inside $E_8$,'' \textit{Commun.\ Math.\ Phys.} \textbf{298}, 419 (2010).

\bibitem{37} J.~C.~Baez, ``Octonions and the Standard Model,'' \url{https://math.ucr.edu/home/baez/standard/} (2020--2025).

\bibitem{38} Springer Nature, ``Artificial intelligence (AI) policy,'' \url{https://www.springernature.com/gp/policies} (2024).

\bibitem{39} J.~A.~Wheeler, ``Information, physics, quantum: the search for links,'' in \textit{Complexity, Entropy, and the Physics of Information}, W.~H.~Zurek (ed.), Addison-Wesley (1990), pp.~3--28.

\bibitem{40} J.~Worrall, ``Structural realism: the best of both worlds?,'' \textit{Dialectica} \textbf{43}, 99--124 (1989).

\bibitem{41} J.~Ladyman, D.~Ross, \textit{Every Thing Must Go: Metaphysics Naturalized}, Oxford University Press (2007).

\bibitem{42} R.~Pin\v{c}\'ak, A.~Pigazzini, M.~Pudl\'ak, E.~Barto\v{s}, ``Geometric origin of a stable black hole remnant from torsion in $G_2$-manifold geometry,'' \textit{Gen.\ Rel.\ Grav.} \textbf{58}, 29 (2026), \href{https://doi.org/10.1007/s10714-026-03528-z}{doi:10.1007/s10714-026-03528-z}.

\bibitem{43} D.~Joyce, S.~Karigiannis, ``A new construction of compact torsion-free $G_2$-manifolds by gluing families of Eguchi--Hanson spaces,'' \textit{J.\ Differential Geom.} (2021), arXiv:1707.09325 (2017).

\bibitem{44} S.~Mukai, ``Finite groups of automorphisms of K3 surfaces and the Mathieu group,'' \textit{Invent.\ Math.} \textbf{94}, 183--221 (1988).

\bibitem{45} A.~Garbagnati, A.~Sarti, ``Symplectic automorphisms of prime order on K3 surfaces,'' \textit{J.\ Algebra} \textbf{318}, 323--350 (2007), arXiv:math/0603742.

\bibitem{46} LEP SUSY Working Group (ALEPH, DELPHI, L3, OPAL), ``Combined LEP chargino results, up to 208 GeV for large $m_0$,'' LEPSUSYWG/01-03.1 (2001), \url{https://lepsusy.web.cern.ch/lepsusy/www/inos_moriond01/charginos_pub.html}.

\bibitem{47} Fermilab, ``DUNE first-beam target of 2031: LBNF/DUNE-US milestone communication'' (2026-05-07), \url{https://news.fnal.gov/}.

\bibitem{48} J.~Chen, H.~Hong, ``Intermediate curvature and splitting theorem,'' arXiv:2604.26529 (2026).

\bibitem{49} M.~Tristram et al., ``Cosmological parameters derived from the final Planck data release (PR4),'' \textit{Astron.\ Astrophys.} \textbf{682}, A37 (2024).

\bibitem{50} E.~Tiesinga, P.~J.~Mohr, D.~B.~Newell, B.~N.~Taylor, ``CODATA recommended values of the fundamental physical constants: 2022,'' \textit{Rev.\ Mod.\ Phys.} \textbf{97}, 025002 (2025), \url{https://physics.nist.gov/cuu/Constants}.

\bibitem{51} V.~V.~Nikulin, ``Integer symmetric bilinear forms and some of their geometric applications,'' \textit{Izv.\ Akad.\ Nauk SSSR Ser.\ Mat.} \textbf{43} (1979); English transl.\ \textit{Math.\ USSR Izv.} \textbf{14} (1980).

\end{thebibliography}
 where the references section should appear.
% Sequential numbering [1]..[50]. Round-1 + round-2 IA-review updates applied
% (2026-07-07): [13] Kovalev 2005, [26] T2K/NOvA Nature 646,
% [27] NuFit 6.0 + 6.1, [45] Garbagnati-Sarti 2007, [46] LEPSUSYWG,
% [47] Fermilab DUNE 2031, [48] Chen-Hong (re-added, cited in S1), [49] Planck PR4, [50] CODATA.

\begin{thebibliography}{99}

\bibitem{1} Particle Data Group, ``Review of Particle Physics,'' \textit{Phys.\ Rev.\ D} \textbf{110}, 030001 (2024).

\bibitem{2} S.~Weinberg, ``Implications of dynamical symmetry breaking,'' \textit{Phys.\ Rev.\ D} \textbf{13}, 974 (1976).

\bibitem{3} Planck Collaboration, ``Planck 2018 results.\ VI.\ Cosmological parameters,'' \textit{Astron.\ Astrophys.} \textbf{641}, A6 (2020).

\bibitem{4} A.~G.~Riess et al., ``A comprehensive measurement of the local value of the Hubble constant,'' \textit{ApJL} \textbf{934}, L7 (2022).

\bibitem{5} C.~D.~Froggatt, H.~B.~Nielsen, ``Hierarchy of quark masses, Cabibbo angles and CP violation,'' \textit{Nucl.\ Phys.\ B} \textbf{147}, 277 (1979).

\bibitem{6} Y.~Koide, ``Fermion-boson two-body model of quarks and leptons and Cabibbo mixing,'' \textit{Lett.\ Nuovo Cim.} \textbf{34}, 201 (1982).

\bibitem{7} C.~Furey, \textit{Standard Model Physics from an Algebra?}, PhD thesis, University of Waterloo (2015).

\bibitem{8} N.~Furey, M.~J.~Hughes, ``One generation of standard model Weyl representations as a single copy of $\mathbb{R} \otimes \mathbb{C} \otimes \mathbb{H} \otimes \mathbb{O}$,'' \textit{Phys.\ Lett.\ B} \textbf{831}, 137186 (2022).

\bibitem{9} R.~Wilson, ``$E_8$ and Standard Model plus gravity,'' arXiv:2404.18938 [hep-th] (2024).

\bibitem{10} T.~P.~Singh et al., ``An $E_8 \otimes E_8$ unification of the Standard Model with pre-gravitation,'' arXiv:2206.06911v3 (2024).

\bibitem{11} B.~S.~Acharya, S.~Gukov, ``M-theory and singularities of exceptional holonomy manifolds,'' \textit{Phys.\ Rep.} \textbf{392}, 121 (2004).

\bibitem{12} L.~Foscolo et al., ``Complete non-compact $G_2$-manifolds from asymptotically conical Calabi--Yau 3-folds,'' \textit{Duke Math.\ J.} \textbf{170}, 3 (2021).

\bibitem{13} A.~Kovalev, ``Coassociative K3 fibrations of compact $G_2$-manifolds,'' arXiv:math/0511150 (2005).

\bibitem{14} M.~Haskins, J.~Nordstr\"om, ``Extra-twisted connected sum $G_2$-manifolds,'' arXiv:1809.09083 (2022); \textit{Ann.\ Glob.\ Anal.\ Geom.} \textbf{64}, art.\ 2 (2023).

\bibitem{15} G.~Bera, ``Associative submanifolds in twisted connected sum $G_2$-manifolds,'' arXiv:2209.00156 (2022).

\bibitem{16} D.~Crowley, S.~Goette, J.~Nordstr\"om, ``An analytic invariant of $G_2$ manifolds,'' \textit{Invent.\ Math.} \textbf{239}(3), 865--907 (2025).

\bibitem{17} T.~Dray, C.~A.~Manogue, \textit{The Geometry of the Octonions}, World Scientific (2015).

\bibitem{18} J.~F.~Adams, \textit{Lectures on Exceptional Lie Groups}, University of Chicago Press (1996).

\bibitem{19} D.~J.~Gross et al., ``Heterotic string theory,'' \textit{Nucl.\ Phys.\ B} \textbf{256}, 253 (1985).

\bibitem{20} D.~D.~Joyce, \textit{Compact Manifolds with Special Holonomy}, Oxford University Press (2000).

\bibitem{21} A.~Kovalev, ``Twisted connected sums and special Riemannian holonomy,'' \textit{J.\ Reine Angew.\ Math.} \textbf{565}, 125 (2003).

\bibitem{22} A.~Corti et al., ``$G_2$-manifolds and associative submanifolds via semi-Fano 3-folds,'' \textit{Duke Math.\ J.} \textbf{164}, 1971 (2015).

\bibitem{23} R.~Harvey, H.~B.~Lawson, ``Calibrated geometries,'' \textit{Acta Math.} \textbf{148}, 47 (1982).

\bibitem{24} B.~S.~Acharya, ``M-theory, $G_2$-manifolds and four-dimensional physics,'' \textit{Class.\ Quant.\ Grav.} \textbf{19}, 5619 (2002).

\bibitem{25} B.~S.~Acharya, E.~Witten, ``Chiral fermions from manifolds of $G_2$ holonomy,'' arXiv:hep-th/0109152 (2001).

\bibitem{26} T2K and NOvA Collaborations, ``Joint analysis of long-baseline neutrino oscillations,'' \textit{Nature} \textbf{646}, 818--824 (2025), \href{https://doi.org/10.1038/s41586-025-09599-3}{doi:10.1038/s41586-025-09599-3}.

\bibitem{27} I.~Esteban et al.\ (NuFIT), ``NuFit-6.0: three-flavour global analysis of neutrino oscillation experiments,'' \textit{JHEP} \textbf{12} (2024) 216, arXiv:2410.05380; NuFIT 6.1 web release, \url{https://www.nu-fit.org} (2025).

\bibitem{28} L.~de Moura, S.~Ullrich, ``The Lean 4 theorem prover and programming language,'' in \textit{CADE 28}, p.\ 625 (2021).

\bibitem{29} mathlib Community, \textit{mathlib4}, \url{https://github.com/leanprover-community/mathlib4}.

\bibitem{30} DUNE Collaboration, ``Long-baseline neutrino facility (LBNF) and DUNE conceptual design report,'' FERMILAB-TM-2696 (2020).

\bibitem{31} DUNE Collaboration, ``Deep Underground Neutrino Experiment (DUNE) far detector technical design report,'' arXiv:2103.04797 (2021).

\bibitem{32} E.~Witten, ``Strong coupling expansion of Calabi-Yau compactification,'' \textit{Nucl.\ Phys.\ B} \textbf{471}, 135 (1996).

\bibitem{33} B.~S.~Acharya et al., ``M-theory solution to the hierarchy problem,'' \textit{Phys.\ Rev.\ D} \textbf{76}, 126010 (2007).

\bibitem{34} M.~Atiyah, E.~Witten, ``M-theory dynamics on a manifold of $G_2$-holonomy,'' \textit{Adv.\ Theor.\ Math.\ Phys.} \textbf{6}, 1 (2002).

\bibitem{35} G.~Kane, \textit{String Theory and the Real World} (2017).

\bibitem{36} J.~Distler, S.~Garibaldi, ``There is no `Theory of Everything' inside $E_8$,'' \textit{Commun.\ Math.\ Phys.} \textbf{298}, 419 (2010).

\bibitem{37} J.~C.~Baez, ``Octonions and the Standard Model,'' \url{https://math.ucr.edu/home/baez/standard/} (2020--2025).

\bibitem{38} Springer Nature, ``Artificial intelligence (AI) policy,'' \url{https://www.springernature.com/gp/policies} (2024).

\bibitem{39} J.~A.~Wheeler, ``Information, physics, quantum: the search for links,'' in \textit{Complexity, Entropy, and the Physics of Information}, W.~H.~Zurek (ed.), Addison-Wesley (1990), pp.~3--28.

\bibitem{40} J.~Worrall, ``Structural realism: the best of both worlds?,'' \textit{Dialectica} \textbf{43}, 99--124 (1989).

\bibitem{41} J.~Ladyman, D.~Ross, \textit{Every Thing Must Go: Metaphysics Naturalized}, Oxford University Press (2007).

\bibitem{42} R.~Pin\v{c}\'ak, A.~Pigazzini, M.~Pudl\'ak, E.~Barto\v{s}, ``Geometric origin of a stable black hole remnant from torsion in $G_2$-manifold geometry,'' \textit{Gen.\ Rel.\ Grav.} \textbf{58}, 29 (2026), \href{https://doi.org/10.1007/s10714-026-03528-z}{doi:10.1007/s10714-026-03528-z}.

\bibitem{43} D.~Joyce, S.~Karigiannis, ``A new construction of compact torsion-free $G_2$-manifolds by gluing families of Eguchi--Hanson spaces,'' \textit{J.\ Differential Geom.} (2021), arXiv:1707.09325 (2017).

\bibitem{44} S.~Mukai, ``Finite groups of automorphisms of K3 surfaces and the Mathieu group,'' \textit{Invent.\ Math.} \textbf{94}, 183--221 (1988).

\bibitem{45} A.~Garbagnati, A.~Sarti, ``Symplectic automorphisms of prime order on K3 surfaces,'' \textit{J.\ Algebra} \textbf{318}, 323--350 (2007), arXiv:math/0603742.

\bibitem{46} LEP SUSY Working Group (ALEPH, DELPHI, L3, OPAL), ``Combined LEP chargino results, up to 208 GeV for large $m_0$,'' LEPSUSYWG/01-03.1 (2001), \url{https://lepsusy.web.cern.ch/lepsusy/www/inos_moriond01/charginos_pub.html}.

\bibitem{47} Fermilab, ``DUNE first-beam target of 2031: LBNF/DUNE-US milestone communication'' (2026-05-07), \url{https://news.fnal.gov/}.

\bibitem{48} J.~Chen, H.~Hong, ``Intermediate curvature and splitting theorem,'' arXiv:2604.26529 (2026).

\bibitem{49} M.~Tristram et al., ``Cosmological parameters derived from the final Planck data release (PR4),'' \textit{Astron.\ Astrophys.} \textbf{682}, A37 (2024).

\bibitem{50} E.~Tiesinga, P.~J.~Mohr, D.~B.~Newell, B.~N.~Taylor, ``CODATA recommended values of the fundamental physical constants: 2022,'' \textit{Rev.\ Mod.\ Phys.} \textbf{97}, 025002 (2025), \url{https://physics.nist.gov/cuu/Constants}.

\bibitem{51} V.~V.~Nikulin, ``Integer symmetric bilinear forms and some of their geometric applications,'' \textit{Izv.\ Akad.\ Nauk SSSR Ser.\ Mat.} \textbf{43} (1979); English transl.\ \textit{Math.\ USSR Izv.} \textbf{14} (1980).

\end{thebibliography}
 where the references section should appear.
% Sequential numbering [1]..[50]. Round-1 + round-2 IA-review updates applied
% (2026-07-07): [13] Kovalev 2005, [26] T2K/NOvA Nature 646,
% [27] NuFit 6.0 + 6.1, [45] Garbagnati-Sarti 2007, [46] LEPSUSYWG,
% [47] Fermilab DUNE 2031, [48] Chen-Hong (re-added, cited in S1), [49] Planck PR4, [50] CODATA.

\begin{thebibliography}{99}

\bibitem{1} Particle Data Group, ``Review of Particle Physics,'' \textit{Phys.\ Rev.\ D} \textbf{110}, 030001 (2024).

\bibitem{2} S.~Weinberg, ``Implications of dynamical symmetry breaking,'' \textit{Phys.\ Rev.\ D} \textbf{13}, 974 (1976).

\bibitem{3} Planck Collaboration, ``Planck 2018 results.\ VI.\ Cosmological parameters,'' \textit{Astron.\ Astrophys.} \textbf{641}, A6 (2020).

\bibitem{4} A.~G.~Riess et al., ``A comprehensive measurement of the local value of the Hubble constant,'' \textit{ApJL} \textbf{934}, L7 (2022).

\bibitem{5} C.~D.~Froggatt, H.~B.~Nielsen, ``Hierarchy of quark masses, Cabibbo angles and CP violation,'' \textit{Nucl.\ Phys.\ B} \textbf{147}, 277 (1979).

\bibitem{6} Y.~Koide, ``Fermion-boson two-body model of quarks and leptons and Cabibbo mixing,'' \textit{Lett.\ Nuovo Cim.} \textbf{34}, 201 (1982).

\bibitem{7} C.~Furey, \textit{Standard Model Physics from an Algebra?}, PhD thesis, University of Waterloo (2015).

\bibitem{8} N.~Furey, M.~J.~Hughes, ``One generation of standard model Weyl representations as a single copy of $\mathbb{R} \otimes \mathbb{C} \otimes \mathbb{H} \otimes \mathbb{O}$,'' \textit{Phys.\ Lett.\ B} \textbf{831}, 137186 (2022).

\bibitem{9} R.~Wilson, ``$E_8$ and Standard Model plus gravity,'' arXiv:2404.18938 [hep-th] (2024).

\bibitem{10} T.~P.~Singh et al., ``An $E_8 \otimes E_8$ unification of the Standard Model with pre-gravitation,'' arXiv:2206.06911v3 (2024).

\bibitem{11} B.~S.~Acharya, S.~Gukov, ``M-theory and singularities of exceptional holonomy manifolds,'' \textit{Phys.\ Rep.} \textbf{392}, 121 (2004).

\bibitem{12} L.~Foscolo et al., ``Complete non-compact $G_2$-manifolds from asymptotically conical Calabi--Yau 3-folds,'' \textit{Duke Math.\ J.} \textbf{170}, 3 (2021).

\bibitem{13} A.~Kovalev, ``Coassociative K3 fibrations of compact $G_2$-manifolds,'' arXiv:math/0511150 (2005).

\bibitem{14} M.~Haskins, J.~Nordstr\"om, ``Extra-twisted connected sum $G_2$-manifolds,'' arXiv:1809.09083 (2022); \textit{Ann.\ Glob.\ Anal.\ Geom.} \textbf{64}, art.\ 2 (2023).

\bibitem{15} G.~Bera, ``Associative submanifolds in twisted connected sum $G_2$-manifolds,'' arXiv:2209.00156 (2022).

\bibitem{16} D.~Crowley, S.~Goette, J.~Nordstr\"om, ``An analytic invariant of $G_2$ manifolds,'' \textit{Invent.\ Math.} \textbf{239}(3), 865--907 (2025).

\bibitem{17} T.~Dray, C.~A.~Manogue, \textit{The Geometry of the Octonions}, World Scientific (2015).

\bibitem{18} J.~F.~Adams, \textit{Lectures on Exceptional Lie Groups}, University of Chicago Press (1996).

\bibitem{19} D.~J.~Gross et al., ``Heterotic string theory,'' \textit{Nucl.\ Phys.\ B} \textbf{256}, 253 (1985).

\bibitem{20} D.~D.~Joyce, \textit{Compact Manifolds with Special Holonomy}, Oxford University Press (2000).

\bibitem{21} A.~Kovalev, ``Twisted connected sums and special Riemannian holonomy,'' \textit{J.\ Reine Angew.\ Math.} \textbf{565}, 125 (2003).

\bibitem{22} A.~Corti et al., ``$G_2$-manifolds and associative submanifolds via semi-Fano 3-folds,'' \textit{Duke Math.\ J.} \textbf{164}, 1971 (2015).

\bibitem{23} R.~Harvey, H.~B.~Lawson, ``Calibrated geometries,'' \textit{Acta Math.} \textbf{148}, 47 (1982).

\bibitem{24} B.~S.~Acharya, ``M-theory, $G_2$-manifolds and four-dimensional physics,'' \textit{Class.\ Quant.\ Grav.} \textbf{19}, 5619 (2002).

\bibitem{25} B.~S.~Acharya, E.~Witten, ``Chiral fermions from manifolds of $G_2$ holonomy,'' arXiv:hep-th/0109152 (2001).

\bibitem{26} T2K and NOvA Collaborations, ``Joint analysis of long-baseline neutrino oscillations,'' \textit{Nature} \textbf{646}, 818--824 (2025), \href{https://doi.org/10.1038/s41586-025-09599-3}{doi:10.1038/s41586-025-09599-3}.

\bibitem{27} I.~Esteban et al.\ (NuFIT), ``NuFit-6.0: three-flavour global analysis of neutrino oscillation experiments,'' \textit{JHEP} \textbf{12} (2024) 216, arXiv:2410.05380; NuFIT 6.1 web release, \url{https://www.nu-fit.org} (2025).

\bibitem{28} L.~de Moura, S.~Ullrich, ``The Lean 4 theorem prover and programming language,'' in \textit{CADE 28}, p.\ 625 (2021).

\bibitem{29} mathlib Community, \textit{mathlib4}, \url{https://github.com/leanprover-community/mathlib4}.

\bibitem{30} DUNE Collaboration, ``Long-baseline neutrino facility (LBNF) and DUNE conceptual design report,'' FERMILAB-TM-2696 (2020).

\bibitem{31} DUNE Collaboration, ``Deep Underground Neutrino Experiment (DUNE) far detector technical design report,'' arXiv:2103.04797 (2021).

\bibitem{32} E.~Witten, ``Strong coupling expansion of Calabi-Yau compactification,'' \textit{Nucl.\ Phys.\ B} \textbf{471}, 135 (1996).

\bibitem{33} B.~S.~Acharya et al., ``M-theory solution to the hierarchy problem,'' \textit{Phys.\ Rev.\ D} \textbf{76}, 126010 (2007).

\bibitem{34} M.~Atiyah, E.~Witten, ``M-theory dynamics on a manifold of $G_2$-holonomy,'' \textit{Adv.\ Theor.\ Math.\ Phys.} \textbf{6}, 1 (2002).

\bibitem{35} G.~Kane, \textit{String Theory and the Real World} (2017).

\bibitem{36} J.~Distler, S.~Garibaldi, ``There is no `Theory of Everything' inside $E_8$,'' \textit{Commun.\ Math.\ Phys.} \textbf{298}, 419 (2010).

\bibitem{37} J.~C.~Baez, ``Octonions and the Standard Model,'' \url{https://math.ucr.edu/home/baez/standard/} (2020--2025).

\bibitem{38} Springer Nature, ``Artificial intelligence (AI) policy,'' \url{https://www.springernature.com/gp/policies} (2024).

\bibitem{39} J.~A.~Wheeler, ``Information, physics, quantum: the search for links,'' in \textit{Complexity, Entropy, and the Physics of Information}, W.~H.~Zurek (ed.), Addison-Wesley (1990), pp.~3--28.

\bibitem{40} J.~Worrall, ``Structural realism: the best of both worlds?,'' \textit{Dialectica} \textbf{43}, 99--124 (1989).

\bibitem{41} J.~Ladyman, D.~Ross, \textit{Every Thing Must Go: Metaphysics Naturalized}, Oxford University Press (2007).

\bibitem{42} R.~Pin\v{c}\'ak, A.~Pigazzini, M.~Pudl\'ak, E.~Barto\v{s}, ``Geometric origin of a stable black hole remnant from torsion in $G_2$-manifold geometry,'' \textit{Gen.\ Rel.\ Grav.} \textbf{58}, 29 (2026), \href{https://doi.org/10.1007/s10714-026-03528-z}{doi:10.1007/s10714-026-03528-z}.

\bibitem{43} D.~Joyce, S.~Karigiannis, ``A new construction of compact torsion-free $G_2$-manifolds by gluing families of Eguchi--Hanson spaces,'' \textit{J.\ Differential Geom.} (2021), arXiv:1707.09325 (2017).

\bibitem{44} S.~Mukai, ``Finite groups of automorphisms of K3 surfaces and the Mathieu group,'' \textit{Invent.\ Math.} \textbf{94}, 183--221 (1988).

\bibitem{45} A.~Garbagnati, A.~Sarti, ``Symplectic automorphisms of prime order on K3 surfaces,'' \textit{J.\ Algebra} \textbf{318}, 323--350 (2007), arXiv:math/0603742.

\bibitem{46} LEP SUSY Working Group (ALEPH, DELPHI, L3, OPAL), ``Combined LEP chargino results, up to 208 GeV for large $m_0$,'' LEPSUSYWG/01-03.1 (2001), \url{https://lepsusy.web.cern.ch/lepsusy/www/inos_moriond01/charginos_pub.html}.

\bibitem{47} Fermilab, ``DUNE first-beam target of 2031: LBNF/DUNE-US milestone communication'' (2026-05-07), \url{https://news.fnal.gov/}.

\bibitem{48} J.~Chen, H.~Hong, ``Intermediate curvature and splitting theorem,'' arXiv:2604.26529 (2026).

\bibitem{49} M.~Tristram et al., ``Cosmological parameters derived from the final Planck data release (PR4),'' \textit{Astron.\ Astrophys.} \textbf{682}, A37 (2024).

\bibitem{50} E.~Tiesinga, P.~J.~Mohr, D.~B.~Newell, B.~N.~Taylor, ``CODATA recommended values of the fundamental physical constants: 2022,'' \textit{Rev.\ Mod.\ Phys.} \textbf{97}, 025002 (2025), \url{https://physics.nist.gov/cuu/Constants}.

\bibitem{51} V.~V.~Nikulin, ``Integer symmetric bilinear forms and some of their geometric applications,'' \textit{Izv.\ Akad.\ Nauk SSSR Ser.\ Mat.} \textbf{43} (1979); English transl.\ \textit{Math.\ USSR Izv.} \textbf{14} (1980).

\end{thebibliography}
